\section{从边郡到州牧：东汉中晚期的中枢失衡}
\label{sec:eastern-han-middle-late-crises}

东汉的危机并不是在一八四年才突然降临。和帝死后，襁褓中的殇帝、安帝以及其后的多次继承，使太后、外戚、宦官和近臣反复成为联结宫廷与官署的枢纽。西羌战争又长期牵动凉州、并州的军费、移民和地方秩序。传世史书擅长记录谁在洛阳获胜，却很少留下边郡军粮如何筹集、流民如何安置的连续账簿；因此不能把一次临朝或一次平乱写成中枢仍能无阻地统治全境的证明。

\subsection{从外戚清算到党锢}

一五九年，桓帝与宦官诛灭梁冀集团，结束了梁氏外戚的公开主导。桓帝个人在政变中的作用、查抄财产的数额和参与者的分工，均有叙事性的夸饰；结果仍然明确：原来依靠外戚维持的宫廷平衡被打破，宦官取得了更直接的进入人事与禁中的渠道。一六六年发生第一次党锢，李膺等官僚、太学生被以结党之名逮捕；一六八年窦武、陈蕃试图诛宦官而失败，次年禁锢扩大。所谓“党人”的名单、地域和持续期限不完全一致，但中枢已经把本可参与官僚评价的士人网络视为敌对的政治资源。

边郡压力没有因此暂停。段颎在一六七年至一六九年间连续进军羌地，战役序列大体可以核对，斩获、迁徙与“平定”的范围却不宜照录。用“羌乱”概括冲突，也不能抹平军民、地方首领和朝廷将领的不同处境。洛阳一面以逮捕和禁锢安排政治忠诚，一面以长期军费、迁徙和暴力处理西部边郡，两个方向都在消耗能让政令被日常执行的关系。

\subsection{石经、学校与被切开的政治世界}

一七五年蔡邕等校定经典并立熹平石经于洛阳太学。残石和太学遗址使立石的事实较为可靠，所刻篇目、完成阶段以及现存石块的归属仍须分开考察。石经把可供比对的文字置于公共空间，显示朝廷仍希望以经典、学校和官员资格组织文化权威；一七八年鸿都门学又招集辞赋、书画、尺牍之士，表明皇帝可另行建立接近宫廷的文才网络。两者不应被写成“儒家”与“非儒家”的简单替换，它们共同提示：人事、知识和声望已经是中枢竞争的一部分。

与此同时，鲜卑边患与汉军出击受挫，使北部防务持续承压。檀石槐所能直接统合的部众、兵数和损失都不能由正史的概称推定。较为确实的是，边境的机动战争与洛阳的党争并行发生，朝廷必须在缺少稳定共识的情况下调集军费、将领和地方合作。

\subsection{黄巾之后，州郡拥有了兵}

一八三年前后，张角以太平道组织信众；次年黄巾起事，冀、豫、兖等地响应，朝廷在各战区征调平叛。起事月份和主要战场较为可靠，参与人数、组织是否完全同步以及各地彼此联络的程度仍有争议。宛城、广宗、下曲阳的黄巾主力先后受挫，并不等于地方社会已经恢复原状；同年凉州军民与羌胡部众又起事，领袖和部众关系屡有变化。朝廷面对的不是一次能够由一场决战结束的叛乱，而是多条道路、地方兵力和征发网络同时失灵。

一八八年改任州牧、州刺史并赋予刘焉等人较大的军政权限，是在这种困局中的应对。制度调整的存在可靠，覆盖范围与州牧实际权力强弱却因州而异，不能把后来的割据格局倒投为当年的统一设计。西园八校尉同年设立，宦官领兵，反映宫廷又想用直接控制的军队平衡外朝将领。中央和地方各自取得武装的渠道，并未形成新的稳定秩序，反而使继承危机更容易转化为调兵。

\subsection{皇帝成为可争夺的资源}

一八九年灵帝死后，何进谋诛宦官而被杀，袁绍等率兵入宫，少帝与陈留王一度被劫出京。屠杀规模、何进召董卓入京的时点和责任分配尚有异说；宫门失守和皇帝脱离宫廷控制的后果却无可回避。董卓随即废少帝、立献帝，并在次年挟天子西迁长安。废立的程序、威逼的细节和太后之死不可虚构成一场无缝的宫廷戏，但皇帝名义自此成为各军必须争夺、又无法由任何一方独占的资源。

关东诸将一九〇年以讨董为名结盟，联盟成员实际参战的程度、盟约约束和兵额都不稳定。袁绍据冀州，曹操迎献帝至许县并以朝廷名义组织屯田，分别显示户口、粮仓、文书与皇帝名义可以被不同军事集团重新组合。二〇〇年官渡之战后，曹操逐步取得河北的仓储与行政核心；乌巢细节、兵力数字和个别来降者的作用不可照传，但袁氏继承内斗与资源失守，使北方的竞争出现了新的中心。

赤壁之战、益州易手、汉中与荆州的争夺，随后把这个中心限定在长江、山地和各地方政权之间。二一九年刘备据汉中称王，关羽北攻樊城后又失荆州；二二〇年曹操死，献帝禅位于曹丕。禅代年月可靠，诏册撰作、群臣参与和“自愿”的性质仍应保留疑问。东汉不是被一场仪式突然替换：州郡军事化、边地战争、宫廷清洗和皇帝名义被各方调用，早已使它失去了能够把这些资源重新合拢的中枢。

\subsection{材料与未被写出的地方}

《后汉书》帝纪、列传和司马彪《续汉书》诸志提供中晚期年序与制度主干，《后汉纪》可用来比对异文；二〇〇年以后的北方战事还必须同《三国志》及其不同政权视角对读。熹平石经残石、边塞与地方考古能校正某些地点和物质条件，却不能替代地方财政、庄园、流民和普通军户的连续记录。本节只据此写出宫廷、边郡、州牧与军队如何互相牵动，不以正史中的人数、族名或人物评语填补基层社会的缺口。
